% 図：ユーザ空間・カーネル空間・ハードウェアの関係（chapters/linux/linux_internals.tex）
\begin{tikzpicture}
  \def\w{9.6}
  % ユーザ空間
  \node[figbox, minimum width=\w cm, minimum height=15mm, fill=green!6, anchor=south west] (user) at (0,4.1) {};
  \node[font=\small\bfseries, anchor=north west] at (user.north west) {ユーザ空間};
  \foreach \t [count=\i from 0] in {シェル（bash）,Web ブラウザ,ROS 2 のノード}
    \node[figbox, minimum width=27mm, minimum height=6mm, font=\footnotesize] at (1.65+3.15*\i, 4.55) {\t};
  % システムコール
  \foreach \x in {1.65,4.8,7.95}
    \draw[figarrow, {Stealth[length=2mm]}-{Stealth[length=2mm]}] (\x,4.1) -- (\x,3.3);
  \node[figlabel, fill=white, inner sep=1pt] at (\w/2,3.7) {システムコール（ファイルを開く、画面に書く、プロセスを作る……）};
  % カーネル空間
  \node[figbox, minimum width=\w cm, minimum height=19mm, fill=blue!10, anchor=south west] (kernel) at (0,1.4) {};
  \node[font=\small\bfseries, anchor=north west] at (kernel.north west) {カーネル空間（Linux カーネル）};
  \foreach \t [count=\i from 0] in {プロセス\\管理,メモリ\\管理,ファイル\\システム,デバイス\\ドライバ,ネット\\ワーク}
    \node[figbox, minimum width=16mm, minimum height=9mm, font=\footnotesize] at (0.95+1.93*\i, 1.95) {\t};
  % ハードウェア
  \draw[figarrow, {Stealth[length=2mm]}-{Stealth[length=2mm]}] (\w/2,1.4) -- (\w/2,0.9);
  \node[figbox, minimum width=\w cm, minimum height=8mm, fill=black!6, anchor=south west] at (0,0)
    {ハードウェア（CPU・メモリ・ストレージ・GPU・入出力装置）};
\end{tikzpicture}
